\section{Introduction}



Large language models (LLMs; \citealp{brown2020language,openai2023gpt4}) have gained increasing popularity as a tool for information seeking.
While they generate engaging and coherent responses,
their outputs are prone to hallucination and often contain factually incorrect information~\cite{ji2023survey}.
This makes it harder for users to trust and verify LLM-generated outputs without any supporting evidence.

\begin{figure}[t]
    \includegraphics[width=0.98\linewidth]{figs/main.pdf}
    \caption{The task setup of \citeeval{}. Given a question, 
    the system generates text while providing \ti{citing passages} from a large retrieval corpus. Each statement may contain multiple citations (e.g., \ttt{[1][2]}). %
    }
    \vspace{-5pt}
    \label{fig:main}
\end{figure}

In this work, we study a new generation paradigm for LLMs, in which
we require LLMs to provide \textit{citations} to one or a few text passages for any statement they generate (Figure~\ref{fig:main}).
Incorporating citations brings several benefits:
(1) users can easily verify LLMs' claims with the provided citations;
(2) LLMs can generate text that faithfully follows cited passages, which has the promise to improve correctness and alleviate hallucination.



\begin{table*}[t]
    \centering
    \small
        \begin{tabular}{lllp{0.5\textwidth}}
            \toprule
            \tf{Dataset} & \tf{Corpus (\#passages)} & \tf{Question type} & \tf{Example} \\
            \midrule
            ASQA & Wikipedia (21M) & Factoid (ambiguous) & \tf{Q}: When did the US break away from England? \\
             &&& \tf{A}: The US declared independence on July 2, 1776 \citedb{[1][2]} ... The Treaty of Paris was later signed on September 3, 1783 \citedb{[3]}. \\ %
            QAMPARI & Wikipedia (21M) & Factoid (list) & \tf{Q}: Which films have Gong Li as a member of their cast? \\
            &&& \tf{A}: The Story of Qiu Ju \citedb{[1]}, Farewell My Concubine \citedb{[2]}, The Monkey King 2 \citedb{[3]}, Mulan \citedb{[3]}, Saturday Fiction \citedb{[3]} ...\\
            ELI5 & Sphere (899M) & Why/How/What & \tf{Q}: How do student loans affect getting a mortgage? \\
            &&& \tf{A}: Student loans can affect the debt to income ratio \citedb{[1]}, which is a key factor in determining the amount that ... \citedb{[2][3]} \\
            \bottomrule
        \end{tabular}
    \caption{The three datasets used in our \citeeval{} benchmark. %
    These datasets cover a wide range of question types
    and the corresponding corpora span from Wikipedia to 
    Web-scale document collection. %
    }
    \vspace{-10pt}
    \label{tab:datasets}
\end{table*}

Multiple commercial systems have adopted this paradigm: %
Bing Chat\footnote{\url{https://www.bing.com/new}} and perplexity.ai\footnote{\url{https://www.perplexity.ai}}
respond to user questions in natural language with references to Web pages.
\citet{nakano2021webgpt,menick2022teaching} share a similar motivation, but they mainly experiment with commercial search engines and closed-source models, making their results difficult to evaluate.
Retrieval-augmented LMs~\citep{borgeaud2022retro,izacard2022atlas}
incorporate retrieved passages during both training and inference,
but do not guarantee faithfulness to retrieved passages or explicitly provide citations.
Additionally,
previous studies mostly rely on human evaluation~\cite{nakano2021webgpt,menick2022teaching,liu2023evaluating}, which is expensive and difficult to reproduce.
We argue that the absence of automated evaluation
hinders the advances of such systems.

We present \tf{\citeeval{}}, the first {reproducible} benchmark for \textit{automatically} evaluating LLMs'  generations with citations.
\citeeval{} assumes a natural-language question and a retrieval corpus, and requires building end-to-end systems to retrieve relevant passages from the corpus, generate a response to the question, and cite corresponding supporting passages.
We compile three datasets that cover different types of questions and corpora---ASQA~\cite{stelmakh2022asqa}, QAMPARI~\cite{rubin2022qampari}, and ELI5~\cite{fan-etal-2019-eli5}---as shown in Table~\ref{tab:datasets}.
Different from previous benchmarks~\citep{lee-etal-2019-latent,bohnet2022attributed},
\citeeval{} evaluates long-text generation, 
focusing on automatically evaluating citation quality,
and allows citing \textit{multiple} passages for individual statements.

We design automatic evaluation methods in three dimensions:
\tf{fluency}, \tf{correctness}, and \tf{citation quality}.
Specifically,
we use MAUVE~\citep{pillutla2021mauve} to measure fluency, 
propose tailored correctness metrics for each dataset, %
and
adopt a natural language inference (NLI) model~\citep{honovich-etal-2022-true-evaluating} to measure citation quality.
We showcase how the three dimensions together contribute to a robust evaluation, 
preventing systems from exploiting shortcuts. %
Additionally, we conduct human evaluation and demonstrate a strong correlation with our automatic metrics.

We experiment on multiple systems with state-of-the-art LLMs and retrievers and also propose novel prompting strategies to synthesize retrieved text into text generation. 
Although all systems are capable of providing fluent and coherent responses, 
there remains substantial room for improvement in terms of correctness and citation quality: For example, on the ELI5 dataset, 
around 50\% generations of our ChatGPT and GPT-4 baselines are not fully supported by the cited passages. 
Additionally, we find that
(1) a closed-book model (generating answers without accessing any retrieved documents) with post-hoc citing achieves good correctness but much worse citation quality;
(2) although interactive retrieval approaches~\citep{yao2023react,schick2023toolformer} offer more flexibility in when/what to retrieve, they do not improve the performance on this challenging benchmark;
(3) summarizing the retrieved passages in a shorter text improves correctness but not citation quality;
(4) reranking multiple generations boosts citation quality measured by human evaluation;
(5) incorporating more retrieved passages in context does not help ChatGPT but improves GPT-4 performance.


Our extensive analyses  
highlight three major challenges of building LLMs to generate text with citations: 
(1) 
the retrieval quality is crucial to the final performance and  has substantial room for improvement;
(2) 
LLMs' limited context window
restricts the number of passages they can incorporate;
(3) 
\revise{current LLMs struggle to synthesize multiple documents in context
without being distracted by irrelevant ones, although better instruction tuning brings significant improvement.}
These challenges pose promising research directions for 
developing better systems integrating retrieval and LLMs.
















